\documentclass{article}
\usepackage[utf8]{inputenc}
\usepackage{amsmath, amsfonts, amssymb} % Para matemática avanzada
\usepackage{todonotes}
\usepackage[margin=2cm]{geometry}

\title{Apuntes de Álgebra I}
\author{Ivan ferryera}
\date{\today}

\begin{document}

\maketitle

\section{ejercicio guia 8}

\textbf{8.} Sea $a, b \in \mathbb{R}$. Probar que para todo $n \in \mathbb{N}$, 

\[
a^n - b^n = (a-b)\sum_{i=1}^{n} a^{i-1}b^{n-i}
\]

Deducir de la formula anterior la serie geometrica : para todo $a \neq 1$,

\[
\sum_{i=1}^{n} a^{i-1} = \frac{a^{n+1}-1}{a-1}
\]

\medskip
\todo[inline]{aca hay que tener cuidado de no marearse, a la hora de leer el ejercicio, ya que nos estan pidiendo probar la primera ecuacion, la segunda ecuacion es una deduccion de la primera, ese sera nuestro segundo problema}
\medskip

Caso base p(1): 

\[
p(1): a^1- b^1 = (a-b)\sum_{i=1}^{1} a^{i-1}b^{1-1}
\]

\[
a- b = (a-b)*1
\]
 
entonces p(1) es verdadero.

\medskip
Paso Inductivo: asumo que que p(k) es verdadero.

\[
p(k): a^k - b^k = (a-b)\sum_{i=1}^{k} a^{i-1}b^{k-i}
\]

Quiero probar que p(k+1) es verdadero:

\[
p(k+1): a^{k+1} - b^{k+1} = (a-b)\sum_{i=1}^{k+1} a^{i-1}b^{k+1-i}
\]

Arranco con el lado derecho de la ecuacion:

\[
(a-b)\sum_{i=1}^{k+1} a^{i-1}b^{k+1-i}  = (a-b)\sum_{i=1}^{k+1} a^{i-1}b^{k-i}b^{1}
\]

\medskip
\todo[inline]{aca hay que separar la sumatoria usando la expresion de p(k), para eso hay que separar el ultimo termino de la sumatoria}
\[
= b.(a-b)\sum_{i=1}^{k+1} a^{i-1}b^{k-i} = b.(a-b)(\sum_{i=1}^{k} a^{i-1}b^{k-i}+ \sum_{i=1}^{k} a^{k}b^{-1})
\]



\end{document}

